\documentclass[10pt,a4paper]{article}
\usepackage[utf8]{inputenc}
\usepackage[vietnamese]{babel}
\usepackage[top=20mm,bottom=20mm,left=18mm,right=18mm,headheight=14pt,headsep=18pt]{geometry}
\usepackage{amsmath,amssymb,amsfonts}
\usepackage{graphicx}
\usepackage{xcolor}
\usepackage{fancyhdr}
\usepackage{newtxtext,newtxmath}
\usepackage{microtype}

% Định nghĩa hệ màu chuẩn sách Stewart Calculus
\definecolor{stewartcyan}{RGB}{0, 130, 185}
\definecolor{stewartgray}{RGB}{100, 100, 100}

% Cấu hình Header và Footer theo chuẩn trang lẻ Stewart
\pagestyle{fancy}
\fancyhf{}
\renewcommand{\headrulewidth}{0pt}
\fancyhead[R]{\parbox{0.71\textwidth}{\textsf{\textbf{MỤC 4.6}}\quad \textsf{Vẽ đồ thị bằng Giải tích và Công nghệ}\hfill\textbf{\textsf{331}}}}
\fancyfoot[C]{\fontsize{7}{9}\selectfont\color{stewartgray} Bản dịch tiếng Việt giáo trình Calculus: Early Transcendentals (9th Edition) -- James Stewart}

\setlength{\parindent}{1.2em}
\setlength{\parskip}{0pt}

\begin{document}

\noindent
\begin{minipage}[t]{0.26\textwidth}
    \vspace{0pt}
    \centering
    % Hình 9
    \includegraphics[width=\linewidth]{images/p0366_fig1.png}
    \vskip 4pt
    {\textsf{\textbf{\footnotesize HÌNH 9}}}

    \vspace{6.6cm}

    % Hình 10
    \includegraphics[width=\linewidth]{images/p0366_fig2.png}
    \vskip 4pt
    {\textsf{\textbf{\footnotesize HÌNH 10}}}

    \vspace{3.2cm}

    % Hình 11
    \includegraphics[width=\linewidth]{images/p0366_fig3.png}
    \vskip 4pt
    {\textsf{\textbf{\footnotesize HÌNH 11}}}
\end{minipage}%
\hfill
\begin{minipage}[t]{0.71\textwidth}
    \vspace{0pt}
    Để ước lượng giá trị cực tiểu, chúng ta phóng to vào hình chữ nhật hiển thị $[-3, 0] \times [-4, 2]$ trong Hình~9. Chúng ta nhận thấy rằng giá trị cực tiểu tuyệt đối xấp xỉ $-3{,}1$ khi $x \approx -0{,}9$, và hàm số nghịch biến trên $(-\infty, -0{,}9)$ cùng $(0, \infty)$, và đồng biến trên $(-0{,}9, 0)$. Các giá trị chính xác nhận được bằng cách lấy đạo hàm:
    \[
        f'(x) = -\frac{7}{x^2} - \frac{6}{x^3} = -\frac{7x + 6}{x^3}
    \]
    Điều này cho thấy $f'(x) > 0$ khi $-\frac{6}{7} < x < 0$ và $f'(x) < 0$ khi $x < -\frac{6}{7}$ và khi $x > 0$. Giá trị cực tiểu chính xác là $f\left(-\frac{6}{7}\right) = -\frac{37}{12} \approx -3{,}08$.

    Hình~9 cũng cho thấy một điểm uốn xuất hiện ở vị trí nào đó giữa $x = -1$ và $x = -2$. Chúng ta có thể ước lượng điểm này chính xác hơn nhiều bằng cách sử dụng đồ thị của đạo hàm cấp hai, nhưng trong trường hợp này, việc tìm các giá trị chính xác cũng dễ dàng tương tự. Vì
    \[
        f''(x) = \frac{14}{x^3} + \frac{18}{x^4} = \frac{2(7x + 9)}{x^4}
    \]
    ta thấy rằng $f''(x) > 0$ khi $x > -\frac{9}{7}$ ($x \neq 0$) và $f''(x) < 0$ khi $x < -\frac{9}{7}$. Do đó $f$ lõm lên trên các khoảng $\left(-\frac{9}{7}, 0\right)$ và $(0, \infty)$, và lõm xuống trên khoảng $\left(-\infty, -\frac{9}{7}\right)$. Điểm uốn là $\left(-\frac{9}{7}, -\frac{71}{27}\right)$.

    Phân tích sử dụng hai đạo hàm đầu tiên cho thấy Hình~8 thể hiện đầy đủ tất cả các khía cạnh quan trọng của đường cong.\hfill\textcolor{stewartcyan}{\rule{1.1ex}{1.1ex}}

    \vspace{1.2em}
    \noindent\textcolor{stewartcyan}{\textbf{\textsf{VÍ DỤ 3}}}\quad Vẽ đồ thị hàm số $f(x) = \dfrac{x^2(x + 1)^3}{(x - 2)^2(x - 4)^4}$.

    \vspace{0.6em}
    \noindent\textcolor{stewartcyan}{\textbf{\textsf{LỜI GIẢI}}}\quad Dựa vào kinh nghiệm với hàm phân thức hữu tỉ trong Ví dụ~2, hãy bắt đầu bằng việc vẽ đồ thị của $f$ trong hình chữ nhật hiển thị $[-10, 10] \times [-10, 10]$. Từ Hình~10, chúng ta có cảm giác rằng mình sẽ phải phóng to để quan sát một số chi tiết tinh vi hơn và cũng phải thu nhỏ để thấy được bức tranh toàn cảnh rộng lớn hơn. Tuy nhiên, như một kim chỉ nam cho việc thu phóng thông minh, trước hết chúng ta hãy xem xét kỹ biểu thức của $f(x)$. Do sự hiện diện của các thừa số $(x - 2)^2$ và $(x - 4)^4$ ở mẫu số, ta dự đoán $x = 2$ và $x = 4$ là các đường tiệm cận đứng. Thật vậy,
    \[
        \lim_{x \to 2} \frac{x^2(x + 1)^3}{(x - 2)^2(x - 4)^4} = \infty \qquad \text{và} \qquad \lim_{x \to 4} \frac{x^2(x + 1)^3}{(x - 2)^2(x - 4)^4} = \infty
    \]
    Để tìm các tiệm cận ngang, chúng ta chia cả tử số và mẫu số cho $x^6$:
    \[
        \frac{x^2(x + 1)^3}{(x - 2)^2(x - 4)^4} = \frac{\dfrac{x^2}{x^3}\,\dfrac{(x + 1)^3}{x^3}}{\dfrac{(x - 2)^2}{x^2}\,\dfrac{(x - 4)^4}{x^4}} = \frac{\dfrac{1}{x}\left(1 + \dfrac{1}{x}\right)^3}{\left(1 - \dfrac{2}{x}\right)^2\left(1 - \dfrac{4}{x}\right)^4}
    \]
    Điều này cho thấy $f(x) \to 0$ khi $x \to \pm\infty$, do đó trục $x$ chính là một tiệm cận ngang.

    Việc khảo sát hành vi của đồ thị gần các giao điểm với trục hoành bằng phân tích tương tự như trong Ví dụ~2.6.12 cũng mang lại hiệu quả rất lớn. Vì $x^2$ luôn dương nên $f(x)$ không đổi dấu tại $0$ và do đó đồ thị của nó không cắt qua trục $x$ tại $0$. Nhưng do có thừa số $(x + 1)^3$, đồ thị có cắt qua trục $x$ tại $-1$ và có một tiếp tuyến nằm ngang tại điểm đó. Kết hợp toàn bộ những thông tin này lại với nhau, mà chưa cần đến đạo hàm, chúng ta thấy rằng đường cong phải có dạng tương tự như đường cong trong Hình~11.
\end{minipage}

\end{document}